\begin{theorem}[Fundamental-group obstruction]\hypertarget{tag-030O}{}\label{push:src:B.4.2}\label{push:B42}
\StatementRecord{push:src:B.4.2}{proved}{030O}{\cite[Theorem 5.1, p.2226]{NO10}}

Let $f:\mathcal Y\to\Delta$ be a
proper holomorphic map from a smooth
complex manifold, with smooth connected
fibres over $\Delta^*$ and reduced
normal-crossings central fibre $D$.
There is a continuous specialization
map from a nearby fibre $Y_t$ to $D$
which induces a surjection
\[
 \pi_1(Y_t)\longrightarrow\pi_1(D).
\]
Hence the central fibre of any such
semistable filling of $S^6$ is simply
connected. In particular it cannot
be a simple pinching of a connected
normalization along one pair of
disjoint nonempty conductor sections,
whether its elliptic normalization
base is split or is a genus-one torsor.
\end{theorem}


\begin{proof}
\begin{enumerate}[label=\arabic*.,leftmargin=*]
\item The surjectivity assertion also follows from
\cite[Theorem~5.1, p.~2226]{NO10}. We use the explicit neighbourhood
retraction needed in \cref{push:src:A.4.7,push:src:B.4.7}. Put
$W=f^{-1}(\overline\Delta_\epsilon)$
and $W^*=W\setminus D$. The negative
gradient of $|f|^2$, for a Hermitian
metric, retracts $W$ onto $D$.
For completeness, in a product chart
$f=z_1\cdots z_k$ the arithmetic--
geometric mean inequality gives
\[
 \|\nabla |f|^2\|
   \ge c\,|f|^{2-1/k}.
\]
A finite cover of compact $D$ gives
a common exponent $\alpha<1$ in
$\|\nabla\rho\|\ge c\rho^\alpha$,
where $\rho=|f|^2$; for a fourfold
one may take $\alpha=7/8$.
The remaining trajectory length is
bounded by
$c^{-1}\rho^{1-\alpha}/(1-\alpha)$.
Hence trajectories converge uniformly
to $D$, and parametrization by the
level of $\rho$ gives a continuous
strong deformation retraction fixing
$D$. General position moves every
loop of $W$ off the real-codimension-two
normal-crossings strata, so
$\pi_1(W^*)\to\pi_1(W)$ is onto.
The smooth proper fibration over the
punctured disc gives
\[
 \pi_1(Y_t)\longrightarrow\pi_1(W^*)
       \longrightarrow\mathbb Z\longrightarrow0.
\]
A small transverse meridian at a
smooth point of a reduced component
maps to the generator of $\mathbb Z$
and bounds a disc in $W$.
Thus the image of $\pi_1(Y_t)$ alone
generates $\pi_1(W)=\pi_1(D)$. Restricting the retraction
$W\to D$ to $Y_t$ gives the asserted continuous specialization
and its surjection on fundamental groups.

\item For the last assertion, let $R$ be
the connected normalization and
$D_0,D_\infty$ its two disjoint
conductor sections. Choose a
continuous function $h:R\to[0,1]$
equal to zero on $D_0$ and one on
$D_\infty$. The function
$\exp(2\pi i h)$ descends to the
pinching. A path from a point of
$D_0$ to its paired point in
$D_\infty$ becomes a loop whose
image has winding number one.
Consequently the pinching has a
surjection of its fundamental group
onto $\mathbb Z$. This contradicts
the preceding necessary condition
when $Y_t\simeq S^6$.
\end{enumerate}
\end{proof}
